\begin{lemma}[ZetaNoZerosInBox]
\label{ZetaNoZerosInBox}
\lean{ZetaNoZerosInBox}
% at /home/runner/work/trough/trough/library/PrimeNumberTheoremAnd/ZetaBounds.lean:3091.0-3214.81
\leanok
For any $T>0$, there is a constant $\sigma<1$ so that
  $$
  \zeta(\sigma'+it) \ne 0
  $$
  for all $|t| \leq T$ and $\sigma' \ge \sigma$.
\end{lemma}
\begin{proof}
\leanok
Assume not. Then there is a sequence $|t_n| \le T$ and $\sigma_n \to 1$ so that
  $\zeta(\sigma_n + it_n) = 0$.
  By compactness, there is a subsequence $t_{n_k} \to t_0$ along which
  $\zeta(\sigma_{n_k} + it_{n_k}) = 0$.
  If $t_0\ne0$, use the continuity of $\zeta$ to get that $\zeta(1 + it_0) = 0$;
  this is a contradiction.
  If $t_0=0$, $\zeta$ blows up near $1$, so can't be zero nearby.
\end{proof}
